\subsection{\texorpdfstring{第~\ref{sec:vector-fields}~节}{向量场}}\label{subsec:solutions-vector-fields}

\begin{description}[font=\normalfont,leftmargin=2em,labelwidth=1.5em]
\RGsolution{ex:prove-lie-bracket-properties}{lie-bracket-properties}
设 $X,Y,Z\in\mathfrak{X}(M)$.
\begin{enumerate}[label=(\alph*),ref=\alph*,leftmargin=*]
\item 设 $\alpha,\beta\in\mathbb{R}$. 则
\begin{align*}
  [\alpha X+\beta Y,Z]&=(\alpha X+\beta Y)\circ Z-Z\circ(\alpha X+\beta Y)\\
  &=\alpha(X\circ Z-Z\circ X)+\beta(Y\circ Z-Z\circ Y)\\
  &=\alpha[X,Z]+\beta[Y,Z].
\end{align*}
类似地, $[X,\alpha Y+\beta Z]=\alpha[X,Y]+\beta[X,Z]$.
\item 有
\[
  [X,Y]=X\circ Y-Y\circ X=-(Y\circ X-X\circ Y)=-[Y,X].
\]
\item Jacobi 恒等式可如下证明:
\begin{align*}
  &[[X,Y],Z]+[[Y,Z],X]+[[Z,X],Y]\\
  &=(X\circ Y-Y\circ X)\circ Z-Z\circ(X\circ Y-Y\circ X)\\
  &\quad +(Y\circ Z-Z\circ Y)\circ X-X\circ(Y\circ Z-Z\circ Y)\\
  &\quad +(Z\circ X-X\circ Z)\circ Y-Y\circ(Z\circ X-X\circ Z)=0.
\end{align*}
\item 设 $f,g\in C^\infty(M)$. 则
\begin{align*}
  [fX,gY]&=(fX)\circ(gY)-(gY)\circ(fX)\\
  &=fgX\circ Y+f(X\cdot g)Y-gfY\circ X-g(Y\cdot f)X\\
  &=fg[X,Y]+f(X\cdot g)Y-g(Y\cdot f)X,
\end{align*}
这里使用了 Exercise~\ref{ex:vector-fields}(\ref{ex:directional-derivative-characterization}).
\end{enumerate}

\RGsolution{ex:rotations-as-vector-field-flows}{rotations-as-vector-field-flows}
\begin{enumerate}[label=(\alph*),ref=\alph*,leftmargin=*]
\item\label{sol:rotation-vector-fields-brackets} 设 $f:\mathbb{R}^3\to\mathbb{R}$ 是光滑函数. 有
\begin{align*}
  [X_1,X_2]\cdot f&=(X_1\circ X_2-X_2\circ X_1)\cdot f\\
  &=X_1\cdot\left(z\frac{\partial f}{\partial x}-x\frac{\partial f}{\partial z}\right)
    -X_2\cdot\left(y\frac{\partial f}{\partial z}-z\frac{\partial f}{\partial y}\right)\\
  &=y\frac{\partial f}{\partial x}-x\frac{\partial f}{\partial y}=-X_3\cdot f,
\end{align*}
因而 $[X_1,X_2]=-X_3$. 类似地, 得到 $[X_1,X_3]=X_2$ 和 $[X_2,X_3]=-X_1$.
\item 设 $V:=\operatorname{span}\{X_1,X_2,X_3\}$. 由 (\ref{sol:rotation-vector-fields-brackets}) 知道, Lie 括号给出双线性映射 $[\cdot,\cdot]:V\times V\to V$, 从而 $V=\operatorname{span}\{X_1,X_2,X_3\}$ 是 $\mathfrak{X}(\mathbb{R}^3)$ 的 Lie 子代数. 为证明它与 $(\mathbb{R}^3,\times)$ 同构, 使用由
\[
  F(a_1X_1+a_2X_2+a_3X_3)=(a_1,-a_2,a_3)
\]
给出的映射 $F:V\to\mathbb{R}^3$. 该映射显然是双射, 因而只需证明它也是 Lie 代数同态. 为此, 注意到
\begin{align*}
  F([X_1,X_2])=F(-X_3)=(0,0,-1)
  &=(1,0,0)\times(0,-1,0)=F(X_1)\times F(X_2);\\
  F([X_1,X_3])=F(X_2)=(0,-1,0)
  &=(1,0,0)\times(0,0,1)=F(X_1)\times F(X_3);\\
  F([X_2,X_3])=F(-X_1)=(-1,0,0)
  &=(0,-1,0)\times(0,0,1)=F(X_2)\times F(X_3).
\end{align*}
\item 设 $q\in\mathbb{R}^3$. 已知 $\psi_{1,t}(q)$ 是 $X_1$ 在 $q$ 处的积分曲线. 因此,
\[
  \frac{\mathrm{d}\psi_{1,t}(q)}{\mathrm{d}t}=(X_1)_{\psi_{1,t}(q)}
  \quad\text{且}\quad\psi_{1,0}(q)=q.
\]
从而
\begin{align*}
  \dot\psi^1_{1,t}(q)&=X_1^1(\psi_{1,t}(q))=0;\\
  \dot\psi^2_{1,t}(q)&=X_1^2(\psi_{1,t}(q))=-\psi^3_{1,t}(q);\\
  \dot\psi^3_{1,t}(q)&=X_1^3(\psi_{1,t}(q))=\psi^2_{1,t}(q).
\end{align*}
故 $\psi^1_{1,t}=C$, $\ddot\psi^2_{1,t}=-\dot\psi^3_{1,t}=-\psi^2_{1,t}$, 因而
\[
  \psi^2_{1,t}=A\cos t+B\sin t
  \quad\text{且}\quad\psi^3_{1,t}=A\sin t-B\cos t,
\]
其中 $A,B,C$ 是 $q=(x,y,z)$ 的函数. 由于 $\psi_{1,0}(x,y,z)=(x,y,z)$, 有 $C=x$, $A=y$ 且 $B=-z$, 从而
\[
  \psi_{1,t}(x,y,z)=(x,y\cos t-z\sin t,y\sin t+z\cos t).
\]
类似地, 有
\[
  \psi_{2,t}(x,y,z)=(x\cos t+z\sin t,y,-x\sin t+z\cos t)
\]
和
\[
  \psi_{3,t}(x,y,z)=(x\cos t-y\sin t,x\sin t+y\cos t,z).
\]
\item 证明 $\psi_{1,\frac{\pi}{2}}\circ\psi_{2,\frac{\pi}{2}}\ne\psi_{2,\frac{\pi}{2}}\circ\psi_{1,\frac{\pi}{2}}$. 其他组合类似. 有
\[
  (\psi_{1,\frac{\pi}{2}}\circ\psi_{2,\frac{\pi}{2}})(x,y,z)
  =\psi_{1,\frac{\pi}{2}}(z,y,-x)=(z,x,y),
\]
而
\[
  (\psi_{2,\frac{\pi}{2}}\circ\psi_{1,\frac{\pi}{2}})(x,y,z)
  =\psi_{2,\frac{\pi}{2}}(x,-z,y)=(y,-z,-x).
\]
\end{enumerate}

\RGsolution{ex:noncomplete-vector-field}{noncomplete-vector-field}
设 $X\in\mathfrak{X}(\mathbb{R})$ 由 $X=x^2\frac{\mathrm{d}}{\mathrm{d}x}$ 给出. 定义其积分曲线的方程为
\[
  \dot x=x^2,
\]
所以在 $x_0\in\mathbb{R}$ 处的积分曲线是 $c_{x_0}(t)=-\frac{1}{t+C}$, 其中 $C$ 是 $x_0$ 的函数. 由于 $c_{x_0}(0)=-\frac{1}{C}$, 得 $C=-\frac{1}{x_0}$, 因而 $X$ 在 $y$ 处的局部流是映射
\[
  \begin{aligned}
  F:W\times(-\varepsilon,\varepsilon)&\to\mathbb{R}\\
  (x,t)&\mapsto\frac{x}{1-tx},
  \end{aligned}
\]
其中 $W=(a,b)$ 是包含 $y$ 的开区间. 例如若 $a>0$, 局部流只能延拓到 $W\times(-\infty,\frac{1}{b})$. 因为局部流无法延拓到 $\mathbb{R}\times\mathbb{R}$, 所以 $X$ 不是完备向量场.

\RGsolution{ex:related-vector-fields-and-integral-curves}{related-vector-fields-and-integral-curves}
\begin{enumerate}[label=(\alph*),ref=\alph*,leftmargin=*]
\item 若 $c:I\to M$ 是 $X$ 的积分曲线, 则对所有 $t\in I$, 有 $\dot c(t)=X_{c(t)}$. 因此, 由 $\gamma(t):=f(c(t))$ 定义的曲线 $\gamma:I\to N$ 满足
\[
  \dot\gamma(t)=(\mathrm{d}f)_{c(t)}\dot c(t)
  =(\mathrm{d}f)_{c(t)}X_{c(t)}=Y_{f(c(t))}=Y_{\gamma(t)},
\]
即 $\gamma$ 是 $Y$ 的积分曲线.
\item 已证得, 若 $X$ 和 $Y$ 关于 $f$ 相关, 则对所有使两边都有定义的 $(p,t)\in M\times\mathbb{R}$, 有 $f(F_X(p,t))=F_Y(f(p),t)$. 反过来, 若此关系成立, 则在 $t=0$ 处求导, 对每个 $p\in M$ 得 $(\mathrm{d}f)_pX_p=Y_{f(p)}$.
\end{enumerate}

\RGsolution{ex:lie-derivative-of-vector-field}{lie-derivative-of-vector-field}
\begin{enumerate}[label=(\alph*),ref=\alph*,leftmargin=*]
\item\label{sol:lie-derivative-is-lie-bracket} 有
\[
  (L_XY)_p=\left.\frac{\mathrm{d}}{\mathrm{d}t}\bigl((\mathrm{d}\psi_{-t})_{\psi_t(p)}Y_{\psi_t(p)}\bigr)\right|_{t=0},
\]
故对任意可微函数 $f\in C^\infty(M)$,
\[
  (L_XY)_p\cdot f=\left.\frac{\mathrm{d}}{\mathrm{d}t}\bigl(Y_{\psi_t(p)}\cdot(f\circ\psi_{-t})\bigr)\right|_{t=0}.
\]
定义
\[
  H(v,\bar v):=Y_{\psi_v(p)}\cdot(f\circ\psi_{-\bar v}).
\]
有
\begin{align*}
  \frac{\partial H}{\partial v}(0,0)
  &=\left.\frac{\mathrm{d}}{\mathrm{d}v}(Y_{\psi_v(p)}\cdot f)\right|_{v=0}
  =\left.\frac{\mathrm{d}}{\mathrm{d}v}(Y\cdot f)(\psi_v(p))\right|_{v=0}\\
  &=(X\cdot(Y\cdot f))(p)
\end{align*}
和
\begin{align*}
  \frac{\partial H}{\partial\bar v}(0,0)
  &=\left.\frac{\mathrm{d}}{\mathrm{d}\bar v}(Y_p\cdot(f\circ\psi_{-\bar v}))\right|_{\bar v=0}
  =Y_p\cdot\left(\left.\frac{\mathrm{d}}{\mathrm{d}\bar v}(f\circ\psi_{-\bar v})\right|_{\bar v=0}\right)\\
  &=Y_p\cdot(-X\cdot f)=-(Y\cdot(X\cdot f))(p).
\end{align*}
因此, 对任意可微函数 $f\in C^\infty(M)$,
\begin{align*}
  (L_XY)_p\cdot f&=\left.\frac{\mathrm{d}}{\mathrm{d}t}(H(t,t))\right|_{t=0}
  =\frac{\partial H}{\partial v}(0,0)+\frac{\partial H}{\partial\bar v}(0,0)\\
  &=(X\cdot(Y\cdot f)-Y\cdot(X\cdot f))(p)=[X,Y]_p\cdot f,
\end{align*}
从而 $L_XY=[X,Y]$.
\item 有
\begin{align*}
  L_X[Y,Z]
  &=\left.\frac{\mathrm{d}}{\mathrm{d}t}((\psi_{-t})_*[Y,Z])\right|_{t=0}
  =\left.\frac{\mathrm{d}}{\mathrm{d}t}[(\psi_{-t})_*Y,(\psi_{-t})_*Z]\right|_{t=0}\\
  &=\left[\left.\frac{\mathrm{d}}{\mathrm{d}t}((\psi_{-t})_*Y)\right|_{t=0},Z\right]
   +\left[Y,\left.\frac{\mathrm{d}}{\mathrm{d}t}((\psi_{-t})_*Z)\right|_{t=0}\right]\\
  &=[L_XY,Z]+[Y,L_XZ].
\end{align*}
注意, 利用 (\ref{sol:lie-derivative-is-lie-bracket}), 这个公式可写成
\[
  [X,[Y,Z]]=[[X,Y],Z]+[Y,[X,Z]]
\]
(即它恰是 Jacobi 恒等式).
\item 有
\[
  L_X(L_YZ)=[X,[Y,Z]]
  \quad\text{且}\quad L_Y(L_XZ)=[Y,[X,Z]].
\]
因此,
\begin{align*}
  (L_X\circ L_Y-L_Y\circ L_X)Z
  &=[X,[Y,Z]]-[Y,[X,Z]]\\
  &=[X,[Y,Z]]+[Y,[Z,X]]\\
  &=-[Z,[X,Y]]=[[X,Y],Z]=L_{[X,Y]}Z,
\end{align*}
其中使用了 Jacobi 恒等式.
\end{enumerate}

\RGsolution{ex:commuting-flows-and-pushforward}{commuting-flows-and-pushforward}
\begin{enumerate}[label=(\alph*),ref=\alph*,leftmargin=*]
\item\label{sol:commuting-flows-from-relatedness} 这是 Exercise~\ref{ex:vector-fields}(\ref{ex:related-vector-fields-and-integral-curves}) 的直接推论.
\item 若对所有 $s,t\in\mathbb{R}$ 有 $\psi_t\circ\phi_s=\phi_s\circ\psi_t$, 则由 (\ref{sol:commuting-flows-from-relatedness}), 对所有 $t\in\mathbb{R}$ 有 $(\psi_t)_*Y=Y$. 因此,
\[
  [X,Y]=L_XY=\left.\frac{\mathrm{d}}{\mathrm{d}t}((\psi_{-t})_*Y)\right|_{t=0}
  =\left.\frac{\mathrm{d}}{\mathrm{d}t}(Y)\right|_{t=0}=0.
\]
另一方面, 若 $[X,Y]=0$, 则
\begin{align*}
  \frac{\mathrm{d}}{\mathrm{d}t}((\psi_t)_*Y)
  &=\left.\frac{\mathrm{d}}{\mathrm{d}\varepsilon}((\psi_{t+\varepsilon})_*Y)\right|_{\varepsilon=0}
  =\left.\frac{\mathrm{d}}{\mathrm{d}\varepsilon}((\psi_t)_*(\psi_\varepsilon)_*Y)\right|_{\varepsilon=0}\\
  &=(\psi_t)_*\left.\frac{\mathrm{d}}{\mathrm{d}\varepsilon}((\psi_\varepsilon)_*Y)\right|_{\varepsilon=0}
  =-(\psi_t)_*L_XY=0.
\end{align*}
由于 $(\psi_0)_*Y=Y$, 得到对所有 $t\in\mathbb{R}$ 有 $(\psi_t)_*Y=Y$. 因此对所有 $s,t\in\mathbb{R}$, $\psi_t\circ\phi_s=\phi_s\circ\psi_t$.
\end{enumerate}
\end{description}
